\chapter{Evaluation \& Experiments}

\section{Experiment 1: Core Data Quality Pipeline}

The main evaluation strategy for DaQuVa is example-driven: each demo script
executes a focused data-quality workflow and writes observable outputs as CSV
or SQLite artifacts. This is appropriate for the prototype stage because the
language is small, the expected behavior is table-oriented, and the produced
tables can be inspected directly.

The existing demo set validates the main pipeline operations:

\begin{itemize}
    \item scanning input tables and annotating rows with tool-generated
          metadata;
    \item filtering by numeric thresholds and string values;
    \item detecting potential duplicates and merging duplicate groups;
    \item exporting cleaned data to CSV files and SQLite databases;
    \item composing several commands into a longer processing pipeline.
\end{itemize}

\section{Experiment 2: Latest Language Features}

The last implementation iteration adds a dedicated showcase demo in
\texttt{code/demo/demo\_showcase/showcase.dqv}. This script exercises the
features added after the initial prototype and acts as a compact five-minute
demonstration scenario.

\begin{lstlisting}[style=dsl]
paid_adults = filter raw where age >= 18 AND plan != "free";
major_providers = filter raw where email contains "@gmail" OR email contains "@outlook";
top5 = sort raw by age desc;
top5 = limit top5 5;
org_users = filter raw where email ends_with ".org";
dupes = find_duplicates adult_paid[name] using tool fuzzy_duplicates(2);
merged = merge dupes;
save merged into clean_customers allowDanger;
\end{lstlisting}

This experiment validates that compound boolean filters respect precedence,
string predicates work on realistic email data, ordering and limiting compose
as expected, duplicate analysis still works after filtering, and the final
cleaned result can be persisted to SQLite only when the script explicitly
contains \texttt{allowDanger}.

\section{Experiment 3: Robustness and Correctness Checks}

Several implementation fixes improve the reliability of the evaluator:

\begin{itemize}
    \item CSV string cells are normalized by stripping surrounding whitespace,
          which makes equality filters behave consistently.
    \item Filters that reference unknown columns now raise an explicit error
          with the available column names.
    \item Ordered comparisons against incompatible non-numeric values now
          produce a clear \texttt{ValueError} instead of an unhandled Python
          exception.
    \item The Levenshtein duplicate matcher short-circuits row pairs whose
          length difference already exceeds the configured distance threshold.
    \item SQLite replacement writes use a temporary table and atomic rename
          pattern, reducing the risk of losing existing data if table creation
          fails.
\end{itemize}

Together, these checks make the prototype more suitable for iterative use:
invalid scripts fail loudly, common data-format issues are handled before
comparison, and persistence operations are more conservative.

\section{Experiment 4: Interactive Use}

The REPL was added to evaluate incremental authoring. Running
\texttt{python -m daquva --repl} starts a persistent session where users can
declare connections, execute one statement at a time, inspect variables with
\texttt{:vars}, inspect active connections with \texttt{:tables}, and print
intermediate tables. The REPL reports row and column counts after assignments,
which makes it easier to validate each transformation before continuing the
pipeline.

\section{Experiment 5: Input/Output Proof Summary}

A data-quality tool should make it easy to prove that processing happened: it
needs a concrete input, observable outputs, and an output that is a
\emph{useful summary} rather than just more annotated columns. The showcase demo
satisfies this end to end. It reads a 30-row input CSV
(\texttt{demo\_showcase/input/customers.csv}), writes cleaned data to both a CSV
file (\texttt{output/cleaned\_customers.csv}) and a SQLite table
(\texttt{clean\_customers}), and finishes with the \texttt{summarize} command,
whose result is exported to \texttt{output/proof\_summary.csv} and persisted to
the SQLite table \texttt{proof\_summary}.

Running the demo produces the proof table in Table~\ref{tab:proof}, which
collapses each pipeline stage to a single row of metrics. It shows directly that
filtering removed five rows (30 raw to 25 paid adults), that duplicate analysis
flagged ten groups across twenty candidate rows, and that merging reduced the
25 candidates to 15 rows while still representing all 25 inputs --- i.e.\ ten
rows were removed by merge.

\begin{table}[H]
\centering
\caption{Aggregate proof summary produced by the showcase demo.}
\label{tab:proof}
\renewcommand{\arraystretch}{1.3}
\begin{tabular}{|l|r|r|r|r|r|}
\hline
\textbf{stage} & \textbf{rows} & \textbf{dup.\ groups} & \textbf{dup.\ cand.} &
\textbf{inputs repr.} & \textbf{removed} \\
\hline
\texttt{raw}         & 30 & 0  & 0  & 30 & 0  \\
\hline
\texttt{paid\_adults} & 25 & 0  & 0  & 25 & 0  \\
\hline
\texttt{dupes}       & 25 & 10 & 20 & 25 & 0  \\
\hline
\texttt{merged}      & 15 & 0  & 0  & 25 & 10 \\
\hline
\end{tabular}
\end{table}

Because the summary is itself a result set, it is verifiable two ways: by
opening the exported CSV, or by querying the \texttt{proof\_summary} table in the
SQLite database alongside the cleaned \texttt{clean\_customers} table. This makes
the input-to-output transformation auditable without re-running the pipeline.
